\chapter{Peluang pada Ruang Terbilang}\label{ch:b2:proba}

Ketiga bab terakhir mengembangkan teori peluang modern: yakni
\omterm{def:b2:proba:space}{ukuran peluang} pada \omterm{def:b2:proba:space}{ruang sampel} yang \omterm{def:b2:structures:countable}{terbilang},
peubah acak diskret, dan \omterm{ex:b2:powerseries:fibonacci}{fungsi pembangkit}. Adapun
teori berhingga pada jilid Kelas 10--12 memperoleh
infrastruktur penuhnya: sebab keaditifan-$\sigma$ menggantikan keaditifan berhingga,
dan perkakas \omterm{def:b2:series:summable}{keluarga terjumlahkan} pada \cref{ch:b2:series}
persis merupakan yang membuat \omterm{def:b2:proba:space}{ruang sampel} tak hingga dapat dikerjakan. Adapun
hasil utamanya di sini adalah kekontinuan peluang sepanjang
barisan \omterm{def:b2:proba:space}{kejadian} yang monoton beserta lema Borel--Cantelli.

\section{Ruang peluang}

\begin{definition}[Ruang peluang terbilang]\label{def:b2:proba:space}
Misalkan $\Omega$ himpunan tak kosong yang berhingga atau \omterm{def:b2:structures:countable}{terbilang} (yakni
\emph{ruang sampelnya}\index{ruang sampel}). Sebuah \emph{ukuran
peluang}\index{ukuran peluang} pada $\Omega$ adalah
pemetaan $\P$ dari himpunan $\mathcal{P}(\Omega)$ berisi semua himpunan bagian
$\Omega$ (yakni \emph{kejadian}\index{kejadian}) ke $[0, 1]$ sedemikian sehingga:
\begin{enumerate}
\item $\P(\Omega) = 1$;
\item (keaditifan-$\sigma$) untuk setiap barisan $(A_n)_{n\in\N}$ berisi
kejadian yang saling lepas berpasangan,
\[
\P\Bigl(\,\bigcup_{n \in \N} A_n\Bigr)
= \sum_{n=0}^{\infty} \P(A_n) .
\]
\end{enumerate}
Adapun pasangan $(\Omega, \P)$ disebut \emph{ruang peluang}\index{ruang
peluang} (yang \omterm{def:b2:structures:countable}{terbilang}).
\end{definition}

\begin{remark}
Pada $\Omega$ yang \omterm{def:b2:structures:countable}{terbilang} kita boleh mengambil semua himpunan bagiannya sebagai \omterm{def:b2:proba:space}{kejadian}; sedangkan pada
ruang yang tak \omterm{def:b2:structures:countable}{terbilang} (sebagaimana diperlukan bagi model \omterm{def:b2:metric:continuity}{kontinu} pada Tahun ke-3)
hal ini tak lagi mungkin, jadi kita membatasi $\P$ pada sebuah
koleksi \omterm{def:b2:proba:space}{kejadian} yang sesuai, yakni sebuah \emph{aljabar-$\sigma$}. Adapun semua rumus
bab ini bertahan pada penyamarataan itu kata demi kata.
\end{remark}

\begin{proposition}[Aturan dasar]\label{prop:b2:proba:rules}
Untuk \omterm{def:b2:proba:space}{kejadian} $A, B$ dan \omterm{def:b2:proba:space}{ukuran peluang} $\P$ berlaku:
$\P(\emptyset) = 0$; $\P$ bersifat aditif berhingga; $\P(A^c) = 1 -
\P(A)$; bila $A \subseteq B$ maka $\P(A) \leq \P(B)$; dan
\[
\P(A \cup B) = \P(A) + \P(B) - \P(A \cap B) .
\]
\end{proposition}

\begin{proof}
Menerapkan keaditifan-$\sigma$ pada $A_0 = \Omega$ dan $A_n = \emptyset$
(untuk $n \geq 1$) memberi $1 = 1 + \sum_{n\geq1}\P(\emptyset)$, jadi
$\P(\emptyset) = 0$; lalu mengganjal gabungan lepas yang berhingga dengan himpunan
kosong memberi keaditifan berhingganya. Adapun sisanya menyusul seperti pada
kasus berhingganya (jilid Kelas 10--12): sebab $1 = \P(A) + \P(A^c)$ dari
$\Omega = A \sqcup A^c$; lalu $\P(B) = \P(A) + \P(B \setminus A)
\geq \P(A)$ ketika $A \subseteq B$; dan setelah diuraikan menjadi tiga
keping yang saling lepas,
\begin{align*}
\P(A \cup B) &= \P(A \setminus B) + \P(B \setminus A) +
\P(A \cap B)\\
&= \bigl(\P(A) - \P(A\cap B)\bigr) + \bigl(\P(B) - \P(A\cap
B)\bigr) + \P(A \cap B),
\end{align*}
yang tak lain inklusi--eksklusi; sedangkan versi $n$ himpunannya yang umum adalah
\cref{exo:b2:proba:4}.
\end{proof}

\begin{proposition}[Distribusi pada ruang
terbilang]\label{prop:b2:proba:distribution}
Memberikan sebuah \omterm{def:b2:proba:space}{ukuran peluang} pada $\Omega =
\{\omega_0, \omega_1, \dots\}$ yang \omterm{def:b2:structures:countable}{terbilang} persis sama dengan memberikan bobot
$p_i = \P(\{\omega_i\}) \geq 0$ dengan $\sum_i p_i = 1$; lalu untuk
setiap $A \subseteq \Omega$,
\[
\P(A) = \sum_{\omega \in A} \P(\{\omega\}) ,
\]
yakni subjumlah (yang konvergen \omterm{def:b2:series:def}{mutlak}) keluarga $(p_i)$.
\end{proposition}

\begin{proof}
Diberikan $\P$, singleton $\{\omega\}$ dengan $\omega \in A$ membentuk
peliput lepas yang \omterm{def:b2:structures:countable}{terbilang} bagi $A$, jadi keaditifan-$\sigma$ memaksa
\[
\P(A) = \sum_{\omega\in A}\P(\{\omega\}),
\]
yakni subjumlah tanpa syarat keluarga taknegatif yang \omterm{def:b2:series:summable}{terjumlahkan}
$(p_i)$ --- sebab penataan ulangnya tak berbahaya persis karena
sukunya taknegatif (\cref{ch:b2:series}); khususnya
$\sum_ip_i = \P(\Omega) = 1$.
Sebaliknya, diberikan bobot taknegatif yang berjumlah total $1$, definisikanlah
$\P(A) = \sum_{\omega \in A}p_\omega$: maka keluarganya \omterm{def:b2:series:summable}{terjumlahkan}, dan
keaditifan-$\sigma$-nya persis merupakan teorema penjumlahan lewat paket
pada \cref{ch:b2:series} yang diterapkan pada pemartisian $\bigcup A_n$
menjadi $A_n$.
\end{proof}

\begin{example}[Model geometri: menunggu gambar
pertama]\label{ex:b2:proba:geometric}
Lemparkan sekeping koin yang berpeluang gambar $p \in \intoo{0}{1}$ berulang kali,
lalu misalkan $\Omega = \N^* \cup \{\infty\}$ mencatat pangkat
gambar pertamanya. Adapun bobot alaminya adalah
\[
\P(\{k\}) = (1 - p)^{k-1}p
\quad (k \in \N^*),
\qquad
\P(\{\infty\}) = 0 ,
\]
yakni sebuah \omterm{def:b2:proba:space}{ukuran peluang} karena $\sum_{k\geq1}(1-p)^{k-1}p =
\frac{p}{1 - (1-p)} = 1$: jadi dengan peluang $1$ permainannya berakhir ---
tetapi \omterm{def:b2:proba:space}{ruang sampelnya} tetap harus memuat kemungkinan bahwa ia
tidak berakhir. Dan keaditifan terbilangnyalah yang memungkinkan kita menegaskan
$\P(\text{permainannya berakhir}) = \sum_k \P(\{k\})$.
\end{example}

\begin{theorem}[Kekontinuan monoton]\label{thm:b2:proba:continuity}
Misalkan $(A_n)$ barisan \omterm{def:b2:proba:space}{kejadian}.
\begin{enumerate}
\item Bila $A_n \subseteq A_{n+1}$ untuk setiap $n$
(yakni \emph{naik}), maka
$\P\bigl(\bigcup_n A_n\bigr) = \lim_{n\to\infty} \P(A_n)$.
\item Bila $A_n \supseteq A_{n+1}$ untuk setiap $n$
(yakni \emph{turun}), maka
$\P\bigl(\bigcap_n A_n\bigr) = \lim_{n\to\infty} \P(A_n)$.
\end{enumerate}
\end{theorem}

\begin{proof}
\emph{1.} Lepaskan: misalkan $B_0 = A_0$ dan $B_n = A_n \setminus
A_{n-1}$. Maka $B_n$ saling lepas berpasangan dengan $\bigcup_{k \leq n}
B_k = A_n$ dan $\bigcup_n B_n = \bigcup_n A_n$. Jadi menurut
keaditifan-$\sigma$ dan keaditifan berhingganya,
\[
\P\Bigl(\bigcup_n A_n\Bigr)
= \sum_{n=0}^\infty \P(B_n)
= \lim_{N\to\infty}\sum_{n=0}^N \P(B_n)
= \lim_{N\to\infty}\P(A_N) .
\]
\emph{2.} Beralihlah ke komplemennya: $(A_n^c)$ naik dengan gabungan
$\bigl(\bigcap A_n\bigr)^c$, lalu terapkanlah bagian 1:
sehingga $1 - \P(\bigcap A_n) = \lim (1 - \P(A_n))$.
\end{proof}

\begin{corollary}[Kesubaditifan terbilang]\label{cor:b2:proba:subadd}
Untuk sembarang barisan \omterm{def:b2:proba:space}{kejadian} berlaku $\P\bigl(\bigcup_n A_n\bigr) \leq
\sum_{n=0}^\infty \P(A_n)$.
\end{corollary}

\begin{proof}
Kesubaditifan berhingganya $\P(A_0 \cup \dots \cup A_N) \leq
\sum_0^N \P(A_n)$ menyusul dari inklusi--eksklusi secara induktif
(atau dari keaditifan atas $B_n \subseteq A_n$ yang telah dilepaskan).
Lalu lewatkan $N \to \infty$: sebab ruas kirinya konvergen ke
$\P(\bigcup_n A_n)$ berkat kekontinuan monoton yang diterapkan pada
barisan naik $C_N = A_0 \cup \dots \cup A_N$.
\end{proof}

\begin{example}[Batas gabungan: kasar tetapi
tak terhancurkan]\label{ex:b2:proba:unionbound}
Kesubaditifan dengan berhingga banyak \omterm{def:b2:proba:space}{kejadian} --- yakni \emph{batas
gabungan} --- menukar ketepatan dengan keuniversalan. Untuk
masalah ulang tahun dengan $23$ orang, membatasi peluang
tabrakannya lewat jumlah atas pasangannya memberi
\[
\P(\text{tabrakan}) \leq \binom{23}2\cdot\frac1{365}
= \frac{253}{365} \approx 0.693 ,
\]
terhadap nilai benarnya $0.507$: jadi meleset dengan selisih lebar, karena
tabrakannya bertumpang tindih. Namun batasnya \emph{tak} menuntut
kebebasan, tak menuntut hukum bersama, dan tak menuntut apa pun selain peluang
pasangannya --- dan itulah sebabnya, pada soal akhir pekan dan
di sepanjang \cref{ch:b2:randomvar}, batas gabungan menjadi
perkakas pertama yang dihunus: sebab ketika ia kebetulan kecil, perkaranya
selesai tanpa pemodelan lebih lanjut.
\end{example}

\begin{example}[Sebuah enam datang, cepat atau lambat]\label{ex:b2:proba:sixeventually}
Lemparkan dadu setimbang selamanya lalu misalkan $B_n = {}$``sekurangnya satu enam
di antara $n$ lemparan pertamanya'', yakni barisan \omterm{def:b2:proba:space}{kejadian} yang naik
dengan $\P(B_n) = 1 - (5/6)^n$. Maka kekontinuan monotonnya memberi
\[
\P(\text{sebuah enam akhirnya muncul})
= \P\Bigl(\bigcup_nB_n\Bigr)
= \lim_n\bigl(1 - (5/6)^n\bigr) = 1 .
\]
Yang penting bukanlah limitnya (yang jelas) melainkan langkah logisnya:
sebab ``akhirnya'' merupakan \omterm{def:b2:proba:space}{kejadian} tentang lemparan yang \emph{tak
berhingga banyaknya}, di luar jangkauan keaditifan berhingga, dan kekontinuan
monoton --- yakni keaditifan-$\sigma$ --- persislah
aksioma yang memberinya sebuah peluang. Jadi setiap pernyataan hampir
pasti pada sisa buku ini melewati pintu sempit yang sama.
\end{example}

\section{Penyaratan dan kebebasan}

\begin{definition}[Peluang bersyarat]\label{def:b2:proba:conditional}
Untuk \omterm{def:b2:proba:space}{kejadian} $A, B$ dengan $\P(B) > 0$, \emph{peluang
bersyarat}\index{peluang bersyarat} bagi $A$ jika diketahui $B$ adalah
\[
\P(A \mid B) = \frac{\P(A \cap B)}{\P(B)} .
\]
Adapun pemetaan $A \mapsto \P(A \mid B)$ sendiri merupakan \omterm{def:b2:proba:space}{ukuran peluang}
pada $\Omega$.
\end{definition}

\begin{remark}
Bahwa $A \mapsto \pcond BA$ kembali merupakan \omterm{def:b2:proba:space}{ukuran peluang}
layak direnungkan sejenak: sebab $\pcond B\Omega = 1$ dan keaditifan-$\sigma$-nya
lolos lewat hasil baginya karena irisan dengan $B$
menghormati gabungan yang lepas. Adapun akibat praktisnya: setiap
kesamaan pada bab ini --- yakni inklusi--eksklusi, kekontinuan
monoton, Borel--Cantelli --- boleh diterapkan \emph{sesudah}
disyaratkan, tanpa bukti baru. Dan itulah sebabnya para ahli peluang terus-menerus
``bekerja di bawah $\pcond B{\cdot}$''.
\end{remark}

\begin{example}[Penyaratan dapat menciptakan
keseragaman]\label{ex:b2:proba:sumseven}
Lemparkan dua dadu setimbang lalu syaratkan jumlahnya $7$: maka untuk
setiap $k \in \intint16$,
\[
\pcond{\{S = 7\}}{X = k}
= \frac{\P(X = k,\ Y = 7 - k)}{\P(S = 7)}
= \frac{1/36}{6/36} = \frac16 :
\]
jadi jika diketahui jumlahnya $7$, dadu pertamanya persis seragam ---
sebab $7$ adalah satu-satunya total yang serasi dengan setiap mukanya, jadi
penyaratannya menghapus semua informasi tentang $X$. Adapun total lain
memiringkan hukumnya (sebab jika diketahui $S = 4$, dadu pertamanya seragam
pada $\{1, 2, 3\}$ saja). Jadi menghitung hukum bersyarat berarti
menormalkan ulang bobot bersamanya sepanjang \omterm{def:b2:proba:space}{kejadian} penyaratannya,
tak lebih.
\end{example}

\begin{example}[Tarikan kedua sama baiknya dengan yang
pertama]\label{ex:b2:proba:seconddraw}
Sebuah guci berisi $3$ bola putih dan $2$ bola hitam; tariklah dua tanpa
pengembalian. Semua orang sepakat $\P(W_1) = \frac35$; lalu berapa
$\P(W_2)$? Peluang total sepanjang tarikan pertamanya:
\[
\P(W_2) = \pcond{W_1}{W_2}\,\P(W_1) +
\pcond{B_1}{W_2}\,\P(B_1)
= \frac24\cdot\frac35 + \frac34\cdot\frac25
= \frac{12}{20} = \frac35 :
\]
yakni persis $\P(W_1)$. Jadi tak ada perhitungan yang diperlukan: sebab berkat kesimetriannya,
setiap bolanya sama mungkin menjadi bola kedua yang ditarik, jadi
tarikan kedua --- \emph{tanpa disyaratkan} --- berhukum sama
dengan yang pertama. Adapun menyaratkan pada hasil pertamanya mengubah
peluangnya; sedangkan tak mengetahuinya tidak. Hujah keterpertukaran
ini kembali pada bab berikutnya bagi pencuplikan tanpa
pengembalian, tempat ia memberi rerata hipergeometrik $np$
tanpa satu kesamaan binomial pun.
\end{example}

\begin{theorem}[Peluang majemuk, peluang total,
Bayes]\label{thm:b2:proba:bayes}
\leavevmode
\begin{enumerate}
\item (Aturan rantai) Bila $\P(A_1 \cap \dots \cap A_{n-1}) > 0$, maka
\[
\P(A_1 \cap \dots \cap A_n)
= \P(A_1)\,\P(A_2 \mid A_1)\cdots
\P(A_n \mid A_1 \cap \dots \cap A_{n-1}) .
\]
\item (Peluang total) Bila $(B_i)_{i \in I}$ merupakan pemartisian $\Omega$
yang berhingga atau \omterm{def:b2:structures:countable}{terbilang} dengan $\P(B_i) > 0$, maka untuk setiap
\omterm{def:b2:proba:space}{kejadian} $A$:
\[
\P(A) = \sum_{i \in I} \P(A \mid B_i)\,\P(B_i) .
\]
\item (Bayes) Di bawah hipotesis yang sama, bila lebih jauh $\P(A) > 0$:
\[
\P(B_j \mid A)
= \frac{\P(A \mid B_j)\,\P(B_j)}
       {\sum_{i \in I} \P(A \mid B_i)\,\P(B_i)} .
\]
\end{enumerate}
\end{theorem}

\begin{proof}
\emph{1.} Tulislah tiap \omterm{def:b2:proba:conditional}{peluang bersyaratnya} sebagai hasil bagi:
maka ruas kanannya adalah
\[
\P(A_1)\cdot\frac{\P(A_1 \cap A_2)}{\P(A_1)}\cdot
\frac{\P(A_1 \cap A_2 \cap A_3)}{\P(A_1 \cap A_2)}\cdots
\frac{\P(A_1 \cap \dots \cap A_n)}{\P(A_1 \cap \dots \cap
A_{n-1})},
\]
yakni hasil kali yang berteleskop: sebab setiap penyebutnya mencoret
pembilang sebelumnya, sehingga menyisakan $\P(A_1 \cap \dots \cap A_n)$.
Adapun semua penyebutnya $\geq \P(A_1 \cap \dots \cap A_{n-1}) >
0$ berkat kemonotonannya, jadi tak ada yang lenyap. (Dan hipotesisnya
menjaga persis ini: sebab menyaratkan pada \omterm{def:b2:proba:space}{kejadian} berpeluang
nol tak terdefinisi.)
\emph{2.} Himpunan $A \cap B_i$ saling lepas berpasangan dengan gabungan
$A$; jadi terapkanlah keaditifan-($\sigma$)-nya beserta definisi
penyaratannya.
\emph{3.} Kedua ruas $\P(B_j \mid A)\P(A) = \P(A \mid
B_j)\P(B_j)$ sama dengan $\P(A \cap B_j)$; jadi bagilah dengan $\P(A)$ lalu uraikan
$\P(A)$ lewat peluang totalnya.
\end{proof}

\begin{example}[Tabrakan ulang tahun, lewat aturan
rantai]\label{ex:b2:proba:birthday}
Dengan $n$ orang yang ulang tahunnya saling bebas dan seragam
atas $365$ hari, misalkan $D_n = {}$``semua $n$ ulang tahunnya berbeda''.
Lalu dengan menyaratkan orang demi orang (yakni aturan rantai):
\[
\P(D_n) = \prod_{k=1}^{n-1}\Bigl(1 - \frac{k}{365}\Bigr),
\]
sebab tiap orang baru harus menghindari $k$ hari yang sudah terpakai.
Untuk $n = 23$: berlaku $\P(D_{23}) \approx 0.493$ --- jadi ulang tahun
yang sama sudah lebih mungkin daripada tidak. Adapun heuristik yang
menjelaskan kekecilan $23$: dengan mengambil logaritmanya, $-\ln
\P(D_n) \approx \sum_{k<n}\frac k{365} =
\frac{\binom n2}{365}$, lalu $\binom{23}2 = 253$ memberi
$253/365 \approx 0.693 \approx \ln 2$. Jadi yang berperan adalah
banyaknya \emph{pasangan}, yang tumbuh kuadratik: sehingga masalah
tabrakan hidup pada skala $n \sim \sqrt{365}$, bukan $n \sim
365$ --- jadi paradoks ulang tahun adalah akar kuadrat yang menyamar.
\end{example}

\begin{example}[Monty Hall, lewat Bayes]\label{ex:b2:proba:montyhall}
Sebuah hadiah bersembunyi di balik satu dari tiga pintu, secara seragam. Kamu memilih
pintu $1$; lalu pembawa acara, yang tahu tempat hadiahnya, membuka salah satu
pintu lainnya yang selalu kosong (dengan memilih seragam ketika ia punya
pilihan), katakanlah pintu $3$. Misalkan $B_i = {}$``hadiahnya di balik pintu
$i$'' dan $A = {}$``pembawa acara membuka pintu $3$''. Maka
$\pcond{B_1}{A} = \frac12$, $\pcond{B_2}{A} = 1$, dan
$\pcond{B_3}{A} = 0$, jadi menurut Bayes
(\cref{thm:b2:proba:bayes}),
\[
\P(B_2 \mid A)
= \frac{1\cdot\frac13}
{\frac12\cdot\frac13 + 1\cdot\frac13 + 0\cdot\frac13}
= \frac23 :
\]
jadi berpindah pintu menang dua kali dari tiga. Adapun perhitungannya
menemukan kekeliruan populernya dengan tepat: sebab langkah pembawa acaranya bersifat
\emph{informatif} (karena ia tak dapat membuka pintu $2$ bila hadiahnya
ada di sana), dan rumus Bayes adalah alat tata buku
yang mengubah ketaksetangkupan ini menjadi $\frac23$. Jadi menyaratkan
pada ``apa yang terlihat'' alih-alih pada ``apa yang benar'' adalah
seluruh seni rumusnya.
\end{example}

\begin{example}[Dua taruhan Chevalier de Méré]\label{ex:b2:proba:demere}
Dua taruhan abad ketujuh belas, yang diselesaikan lewat kebebasan. Taruhan
pertama: sekurangnya satu enam dalam $4$ lemparan sebuah dadu,
\[
\P = 1 - \Bigl(\frac56\Bigr)^{\!4} \approx 0.518 > \frac12 .
\]
Taruhan kedua: sekurangnya satu enam ganda dalam $24$ lemparan dengan dua dadu,
\[
\P = 1 - \Bigl(\frac{35}{36}\Bigr)^{\!24} \approx 0.491 <
\frac12 .
\]
De Méré menalar bahwa $24$ lemparan berpeluang $\frac1{36}$
sepatutnya menyamai $4$ lemparan berpeluang $\frac16$ (sebab nisbahnya sama
$\frac{24}{36} = \frac46$); dan kegagalan kesebandingan
ini --- karena peluang gabungan tak menskala
secara linear --- konon mendorong suratnya kepada Pascal,
sehingga melahirkan teori peluang. Adapun pembandingan
yang benar lewat logaritma: sebab $n$ percobaan berpeluang $p$
berhasil sekurangnya sekali dengan peluang $1 - (1-p)^n \approx
1 - \eu^{-np}$, jadi invarian yang jujurnya adalah $np$: di sini
$4\cdot\frac16 = \frac23$ lawan $24\cdot\frac1{36} =
\frac23$ --- jadi sama! Sehingga kedua taruhannya hanya berselisih pada orde
kedua dalam $p$, dan cukup untuk memindahkan satu di antaranya menyeberangi
garis lima puluh persen: jadi peluang kecil adalah ranah tempat
gerak hati memerlukan eksponensial, bukan penggaris.
\end{example}

\begin{remark}[Kekeliruan yang sering muncul pada penyaratan]
Ada tiga kerancuan yang berulang, yang semuanya terlihat pada contoh
di atas. (i) \emph{Pembalikan}: sebab $\pcond BA$ dan $\pcond AB$
berselisih faktor $\P(A)/\P(B)$ --- jadi sebuah uji yang
$99\%$ akurat pada yang sakit masih boleh menyisakan pasien yang positif
hampir pasti sehat ketika penyakitnya langka
(\cref{exo:b2:proba:3}); sehingga mengutip $\pcond{\text{sakit}}{
\text{positif}}$ ketika yang dimaksud $\pcond{\text{positif}}{
\text{sakit}}$ adalah kekeliruan laju-dasarnya. (ii)
\emph{Menyaratkan pada \omterm{def:b2:proba:space}{kejadian} yang keliru}: sebab pada Monty Hall,
\omterm{def:b2:proba:space}{kejadian} penyaratan yang benar adalah ``pembawa acara membuka pintu $3$'',
bukan ``hadiahnya tak di balik pintu $3$''; karena keduanya mengusung
informasi yang berbeda, dan seluruh $\frac23$-nya bergantung pada
selisih itu. (iii) \emph{Saling lepas lawan \omterm{def:b2:proba:independence}{saling bebas}}:
sebab \omterm{def:b2:proba:space}{kejadian} yang saling lepas dan berpeluang positif tak pernah
\omterm{def:b2:proba:independence}{saling bebas} ($\P(A\cap B) = 0 \neq \P(A)\P(B)$) ---
jadi kebebasan adalah keserasian informasi, bukan ketiadaan
tumpang tindih.
\end{remark}

\begin{definition}[Kebebasan]\label{def:b2:proba:independence}
\omterm{def:b2:proba:space}{Kejadian} $A$ dan $B$ disebut \emph{saling bebas}\index{kejadian saling bebas} bila $\P(A \cap B) =
\P(A)\P(B)$. Adapun keluarga \omterm{def:b2:proba:space}{kejadian} $(A_i)_{i \in I}$ disebut
\emph{saling bebas (secara bersama)} bila untuk setiap himpunan bagian berhingga $J
\subseteq I$,
\[
\P\Bigl(\bigcap_{i \in J} A_i\Bigr)
= \prod_{i \in J} \P(A_i) .
\]
\end{definition}

\begin{remark}
Kebebasan bersama tegas lebih kuat daripada kebebasan
berpasangan: sebab dengan dua lemparan koin setimbang, \omterm{def:b2:proba:space}{kejadian} ``yang pertama
gambar'', ``yang kedua gambar'', dan ``keduanya sepakat'' \omterm{def:b2:proba:independence}{saling bebas}
berpasangan (sebab tiap pasangannya beririsan dengan peluang $\frac14 =
\frac12\cdot\frac12$), namun irisan bertiganya berpeluang
$\frac14 \neq \frac18$. Perhatikan pula bahwa bila $A, B$ \omterm{def:b2:proba:independence}{saling bebas},
maka $A, B^c$ juga (hitunglah: $\P(A \cap B^c) = \P(A) - \P(A\cap B) =
\P(A)(1 - \P(B))$), sehingga $A^c, B^c$ juga.
\end{remark}

\begin{example}[Kebebasan yang terbaca dari struktur hasil
kali]\label{ex:b2:proba:productindep}
Lemparkan dua dadu setimbang: $\Omega = \intint16^2$ dengan bobot
seragam. Misalkan $A = {}$``dadu pertama genap'' dan $B = {}$``dadu
kedua sekurangnya $5$''. Setelah dicacah: $\abs A = 3\cdot6 = 18$,
$\abs B = 6\cdot2 = 12$, $\abs{A\cap B} = 3\cdot2 = 6$, jadi
\[
\P(A\cap B) = \frac6{36} = \frac{18}{36}\cdot\frac{12}{36} =
\P(A)\,\P(B) :
\]
jadi \omterm{def:b2:proba:independence}{saling bebas}, dan mekanismenya kasatmata --- sebab $A$ mengendalai
koordinat pertamanya saja, $B$ koordinat keduanya saja, dan
ukuran seragam pada himpunan hasil kali membuat cacah koordinatnya
mengalikan. Jadi setiap klaim bertipe ``\omterm{def:b2:proba:space}{kejadian} yang bergantung pada
kelompok lemparan yang saling lepas bersifat \omterm{def:b2:proba:independence}{saling bebas}'' (yang dipakai besar-besaran
pada soal akhir pekan) adalah perhitungan ini, dengan indeks yang lebih banyak.
\end{example}

\begin{example}[Analisis langkah pertama]\label{ex:b2:proba:firststep}
Untuk model geometri pada \cref{ex:b2:proba:geometric},
berapa peluang $u$ bahwa gambar pertamanya jatuh pada pangkat yang
\emph{genap}? Syaratkanlah pada lemparan pertamanya: dengan peluang
$p$ pangkatnya $1$ (yang ganjil); sedangkan dengan peluang $q = 1 - p$
permainannya bermula ulang dengan semua paritasnya terbalik, jadi
\[
u = p\cdot0 + q\,(1 - u)
\qquad\Longrightarrow\qquad
u = \frac{q}{1 + q} .
\]
Jadi satu baris, tanpa deret --- dan ia sepakat dengan penjumlahan
langsung pada \cref{exo:b2:proba:9}, yang memberi $1 - u =
\frac1{1+q}$. Adapun teknik ``langkah pertama'' ini (yakni menyaratkan pada
percobaan pertamanya lalu mengenali salinan masalahnya yang tergeser) adalah
bentuk peluang bagi sebuah rekursi, dan ia
mesin di balik persamaan lama permainan pada
\cref{exo:b2:proba:6} beserta perhitungan singgahan pertama pada
soal akhir pekan.
\end{example}

\section{Lema Borel--Cantelli}

\begin{definition}[Limit superior kejadian]\label{def:b2:proba:limsup}
Untuk sebuah barisan \omterm{def:b2:proba:space}{kejadian} $(A_n)$, \omterm{def:b2:proba:space}{kejadian}
\[
\limsup_n A_n
= \bigcap_{N=0}^{\infty}\ \bigcup_{n \geq N} A_n
= \{\omega \in \Omega : \omega \in A_n
\text{untuk tak berhingga banyak } n\}
\]
adalah \omterm{def:b2:proba:space}{kejadian} ``$A_n$ terjadi tak berhingga sering''.
\end{definition}

\begin{example}[Menerjemahkan ``tak berhingga sering'' dan
``akhirnya'']\label{ex:b2:proba:limsuplang}
Komplemen $\limsup_nA_n$ adalah, menurut de Morgan,
\[
\Bigl(\bigcap_N\bigcup_{n\geq N}A_n\Bigr)^{\!c}
= \bigcup_N\bigcap_{n\geq N}A_n^c
= \{\omega : \omega \notin A_n \text{ untuk setiap } n \text{ yang besar}\},
\]
yakni \omterm{def:b2:proba:space}{kejadian} ``\emph{akhirnya}, $A_n$ gagal'' (yang ditulis
$\liminf_nA_n^c$). Jadi ``$A_n$ tak berhingga sering'' dan
``$A_n^c$ akhirnya'' saling berkomplemen --- dan menjaga
kamus ini tetap lurus mencegah sebagian besar kecelakaan kuantor.
Beberapa terjemahan contoh bagi pelemparan koin: ``gambar yang
tak berhingga banyaknya'' adalah $\limsup\{X_n = H\}$; ``hanya berhingga banyak deretan
$100$ gambar'' adalah komplemen sebuah limsup; dan ``frekuensi
berjalannya konvergen ke $\frac12$'' adalah
$\bigcap_j\bigcup_N\bigcap_{n\geq N}\{\abs{\widehat p_n -
\tfrac12} < \tfrac1j\}$ --- jadi operasi \omterm{def:b2:structures:countable}{terbilang} di seluruhnya,
sehingga semuanya \omterm{def:b2:proba:space}{kejadian} yang jujur.
\end{example}

\begin{theorem}[Borel--Cantelli]\label{thm:b2:proba:borelcantelli}
\leavevmode
\begin{enumerate}
\item Bila $\sum_{n} \P(A_n) < \infty$, maka
$\P\bigl(\limsup_n A_n\bigr) = 0$.
\item Bila \omterm{def:b2:proba:space}{kejadian} $A_n$ \omterm{def:b2:proba:independence}{saling bebas} dan $\sum_n \P(A_n) =
\infty$, maka $\P\bigl(\limsup_n A_n\bigr) = 1$.
\end{enumerate}
\end{theorem}

\begin{proof}
\emph{1.} Misalkan $C_N = \bigcup_{n \geq N}A_n$; maka barisan $(C_N)$
turun dengan irisan $\limsup A_n$, jadi menurut kesubaditifan
\omterm{def:b2:structures:countable}{terbilangnya} (\cref{cor:b2:proba:subadd})
\[
\P(C_N) \leq \sum_{n \geq N}\P(A_n)
\xrightarrow[N\to\infty]{} 0
\]
(yakni ekor sebuah deret yang konvergen). Lalu kekontinuan monoton
(\cref{thm:b2:proba:continuity}) merampungkannya: $\P(\limsup A_n) =
\lim_N \P(C_N) = 0$.

\emph{2.} Cukuplah menunjukkan $\P\bigl(\bigcup_{n\geq N}A_n\bigr)
= 1$ untuk setiap $N$: sebab bila \omterm{def:b2:proba:space}{kejadian} $B_N$ semuanya berpeluang
$1$, maka
\[
\P\Bigl(\Bigl(\bigcap_NB_N\Bigr)^{\!c}\Bigr)
= \P\Bigl(\bigcup_NB_N^c\Bigr)
\leq \sum_N\P(B_N^c) = 0
\]
menurut kesubaditifan \omterm{def:b2:structures:countable}{terbilangnya}
(\cref{cor:b2:proba:subadd}), jadi irisan \omterm{def:b2:structures:countable}{terbilang}
$\limsup A_n = \bigcap_N\bigcup_{n\geq N}A_n$ tetap berpeluang
$1$. Tetapkan $N$, lalu tinjaulah untuk $M > N$
komplemennya:
\[
\P\Bigl(\bigcap_{n=N}^{M} A_n^c\Bigr)
= \prod_{n=N}^{M}\bigl(1 - \P(A_n)\bigr)
\leq \prod_{n=N}^{M} e^{-\P(A_n)}
= \exp\Bigl(-\sum_{n=N}^M \P(A_n)\Bigr) ,
\]
dengan memakai kebebasan komplemennya beserta batas kecembungan $1 -
x \leq e^{-x}$. Lalu ketika $M \to \infty$ eksponennya menuju $-\infty$
berkat kedivergenan deretnya, jadi menurut kekontinuan monoton
(bagi barisan yang turun) $\P\bigl(\bigcap_{n \geq N}A_n^c\bigr) = 0$,
yakni $\P\bigl(\bigcup_{n \geq N}A_n\bigr) = 1$.
\end{proof}

\begin{example}[Deretan gambar yang tak berhingga]\label{ex:b2:proba:runs}
Lemparkan koin setimbang selamanya, lalu misalkan $A_n$ \omterm{def:b2:proba:space}{kejadian} ``lemparan $n,
n+1, \dots, n + k - 1$ semuanya gambar'' (yakni deretan $k$ gambar
yang bermula pada waktu $n$), untuk $k$ yang tetap. Maka \omterm{def:b2:proba:space}{kejadian} $A_{jk}$ (dengan $j =
1, 2, \dots$), yang bergantung pada blok lemparan yang saling lepas, bersifat
\omterm{def:b2:proba:independence}{saling bebas}, masing-masing berpeluang $2^{-k}$, dan $\sum_j 2^{-k} =
\infty$: jadi menurut Borel--Cantelli~2, dengan peluang $1$ tak berhingga
banyak bloknya seluruhnya gambar --- sehingga \emph{setiap} pola yang tetap berulang
tak berhingga sering, secara hampir pasti. Sebaliknya, bila \omterm{def:b2:curves:length}{panjang} deretannya
kita biarkan tumbuh, $B_n = {}$``sebuah deretan $2\log_2 n$ gambar bermula di
$n$'' punya $\P(B_n) = n^{-2}$ yang \omterm{def:b2:series:summable}{terjumlahkan}, jadi hampir pasti hanya
berhingga banyak deretan \omterm{def:b2:curves:length}{panjang} semacam itu bermula: sehingga Borel--Cantelli menakar
persis \emph{seberapa \omterm{def:b2:curves:length}{panjang}} deretan terpanjangnya.
\end{example}

\begin{example}[Monyet tak hingga, dikuantifikasi]
Seekor monyet mengetik huruf seragam yang \omterm{def:b2:proba:independence}{saling bebas} dari abjad
berhuruf $26$. Potonglah ketikannya menjadi blok lepas
berisi empat huruf; maka \omterm{def:b2:proba:space}{kejadian} $A_j = {}$``blok $j$ mengeja
MATH'' \omterm{def:b2:proba:independence}{saling bebas} dengan $\P(A_j) = 26^{-4}$, dan
$\sum_j\P(A_j) = \infty$: jadi menurut Borel--Cantelli~2 monyetnya
mengetik MATH tak berhingga sering, secara hampir pasti --- dan hal yang sama
berlaku bagi sembarang teks tetap sepanjang apa pun, dengan bloknya disesuaikan.
Adapun catatan kuantitatifnya mengempiskan keajaibannya: sebab $26^4 =
456\,976$, jadi MATH pertamanya menuntut sekitar setengah juta
ketukan tuts secara rata-rata, dan sebuah lakon Shakespeare berisi $10^5$
karakter menunggu berorde $26^{10^5}$ blok --- jadi hampir pasti
adalah pernyataan tentang cakrawala $\infty$, bukan tentang cakrawala
mana pun yang akan ditemui seekor monyet. Jadi Borel--Cantelli mengesahkan
limitnya; sedangkan ukuran sukunya yang menceritakan pada skala manusia.
\end{example}

\begin{remark}
Pada \cref{ex:b2:proba:runs} \omterm{def:b2:proba:space}{ruang sampel} yang mendasarinya (yakni barisan
lemparan yang tak hingga) bersifat tak \omterm{def:b2:structures:countable}{terbilang}, jadi secara ketat
contohnya hidup pada kerangka teori ukuran Tahun ke-3; namun
\emph{perhitungannya} hanya memakai aturan yang dibuktikan pada bab
ini, yang diterapkan pada \omterm{def:b2:proba:space}{kejadian} yang ditentukan berhingga banyak lemparan dan
kombinasi \omterm{def:b2:structures:countable}{terbilangnya}. Inilah konvensi baku pada tingkat ini:
teorinya dinyatakan pada ruang \omterm{def:b2:structures:countable}{terbilang}, sedangkan contoh permainan
tak hingga ditangani dengan perkakas yang sama.
\end{remark}

\begin{remark}[Pandangan ke depan di dalam jilid ini]
Perkakas bab ini dikonsumsi borongan oleh kedua bab
berikutnya. Indikator mengubah \omterm{def:b2:proba:space}{kejadian} menjadi peubah acak, dan
keaditifan-$\sigma$ menjadi keterjumlahan yang mendefinisikan
nilai harapan (\cref{ch:b2:randomvar}); sedangkan Borel--Cantelli ditambah
batas ekor yang \omterm{def:b2:series:summable}{terjumlahkan} persis merupakan cara hukum bilangan besar
yang kuat bagi koin dibuktikan di sana. Adapun pada
\cref{ch:b2:genfun}, kekontinuan monotonnya muncul kembali pada saat
yang menentukan: sebab peluang kepunahan sebuah proses bercabang
\emph{didefinisikan} sebagai limit monoton
$\lim\P(Z_n = 0)$, dan persamaan titik tetap yang dipenuhinya
diperoleh dengan melewatkan limit pada barisan naik itu ---
jadi teorema terakhir buku ini berdiri di atas teorema pertama bab ini.
\end{remark}

\begin{remark}[Metode: tiga jalan menuju peluang satu]
Pernyataan hampir pasti dibuktikan lewat tiga tuas, yakni dalam
urutan kekuatan yang menaik. \emph{Kekontinuan monoton}:
pamerkanlah \omterm{def:b2:proba:space}{kejadiannya} sebagai gabungan yang naik (atau irisan
yang turun) atas \omterm{def:b2:proba:space}{kejadian} bercakrawala berhingga yang peluangnya
terhitung (\cref{ex:b2:proba:sixeventually}).
\emph{Gabungan nol}: sebab gabungan \omterm{def:b2:structures:countable}{terbilang} atas \omterm{def:b2:proba:space}{kejadian} berpeluang nol
bernilai nol (berkat kesubaditifan \omterm{def:b2:structures:countable}{terbilangnya}), jadi cukuplah
membunuh tiap \omterm{def:b2:proba:space}{kejadian} buruknya secara terpisah --- dan begitulah ``untuk setiap
$j$, akhirnya $\abs{\widehat p_n - p} < 1/j$'' terangkai
menjadi kekonvergenan. \emph{Borel--Cantelli}: ketika \omterm{def:b2:proba:space}{kejadiannya} berupa
limsup, jumlahkanlah peluangnya; sebab kekonvergenannya membunuhnya
(tanpa perlu kebebasan), sedangkan kedivergenan ditambah kebebasan
mengesahkannya. Jadi memilih tuas yang tepat biasanya seluruh
buktinya; dan soal akhir pekan menjalankan ketiganya dalam satu
hujah.
\end{remark}

\begin{remark}[Di mana ini dipakai]
Kekontinuan monoton dan Borel--Cantelli adalah dua tuas bagi
setiap pernyataan ``hampir pasti'': sebab keduanya menggerakkan kerekurenan
\omterm{pb:b2:proba:1}{jalan acak} pada soal akhir pekan bab ini,
sisi hampir pasti hukum bilangan besar
(\cref{ch:b2:randomvar}), dan analisis kepunahan
proses bercabang (\cref{ch:b2:genfun}). Adapun jilid Tahun ke-3
membangun ulang teorinya atas aljabar-$\sigma$ dan pengintegralan
Lebesgue, tempat \omterm{def:b2:proba:space}{ruang sampel} tak \omterm{def:b2:structures:countable}{terbilang} yang dipakai
secara tak resmi di sini menjadi sepenuhnya ketat.
\end{remark}

\section{Latihan}

\begin{exercise}[$\star$]\label{exo:b2:proba:1}
Sebuah guci berisi $n$ bola bernomor. Bolanya ditarik satu demi satu
tanpa pengembalian. Hitunglah peluang bahwa bola bernomor $1$
tertarik sebelum bola bernomor $2$. Samaratakan: yakni peluang bahwa
bola $1$ tertarik pertama di antara bola $1, \dots, k$.
\end{exercise}

\begin{exercise}[$\star$]\label{exo:b2:proba:2}
Tunjukkan bahwa pada $\Omega = \N^*$ bobot $p_k = \frac{1}{k(k+1)}$
mendefinisikan sebuah \omterm{def:b2:proba:space}{ukuran peluang}, lalu hitunglah $\P(2\N^*)$ (yakni hasil
yang genap) sebagai sebuah deret; kemudian tunjukkan bahwa ia sama dengan $1 - \ln 2$.
\emph{(Teleskopkan $\frac{1}{2j(2j+1)} = \frac{1}{2j} -
\frac{1}{2j+1}$ lalu pakailah deret harmonik yang berselang-seling,
\cref{ch:b2:series}.)}
\end{exercise}

\begin{exercise}[$\star$]\label{exo:b2:proba:3}
(Positif palsu) Sebuah penyakit menjangkiti satu orang dari $10\,000$. Sebuah
uji mendeteksinya dengan peluang $0.99$ pada yang sakit, dan memberi
positif palsu dengan peluang $0.01$ pada yang sehat. Hitunglah
peluang menjadi sakit jika diketahui ujinya positif, lalu berilah komentar.
\end{exercise}

\begin{exercise}[$\star\star$]\label{exo:b2:proba:4}
Misalkan $A_1, \dots, A_n$ \omterm{def:b2:proba:space}{kejadian}. Buktikan rumus
inklusi--eksklusi
\[
\P\Bigl(\bigcup_{i=1}^n A_i\Bigr)
= \sum_{\emptyset \neq J \subseteq \{1,\dots,n\}}
(-1)^{\abs J + 1}\,\P\Bigl(\bigcap_{i \in J}A_i\Bigr)
\]
dengan mengintegralkan kesamaan $1 - \prod_{i=1}^n(1 -
\mathbf{1}_{A_i}) = \mathbf{1}_{\bigcup A_i}$ atas $\Omega$
(yakni menjumlahkannya terbobot oleh $\P(\{\omega\})$).
\end{exercise}

\begin{exercise}[$\star\star$]\label{exo:b2:proba:5}
(Masalah pencocokan, lewat inklusi--eksklusi) Sebanyak $n$ surat dimasukkan
secara acak seragam ke $n$ amplop, satu per amplop. Dengan memakai
\cref{exo:b2:proba:4}, tunjukkanlah bahwa peluang \emph{tak ada}
kecocokan adalah $\sum_{k=0}^n \frac{(-1)^k}{k!} \to e^{-1}$, lalu
turunkan peluang tepat satu kecocokan.
\end{exercise}

\begin{exercise}[$\star\star$]\label{exo:b2:proba:6}
Sekeping koin yang berat sebelah (dengan peluang gambar $p \in \intoo{0}{1}$) dilempar
sampai dua gambar berurutan muncul. Misalkan $q_n$ peluang
bahwa permainannya berlangsung lebih dari $n$ lemparan. Tunjukkan, dengan menyaratkan pada
lemparan pertamanya, bahwa $q_n = (1-p)\,q_{n-1} + p(1-p)\,q_{n-2}$
untuk $n \geq 2$, lalu turunkan bahwa permainannya berakhir dengan peluang
$1$. \emph{(Tunjukkan $q_n \to 0$ dengan membandingkannya dengan barisan
geometri: sebab kedua akar persamaan karakteristiknya bernilai mutlak di
$\intoo{0}{1}$.)}
\end{exercise}

\begin{exercise}[$\star\star\star$]\label{exo:b2:proba:7}
(Rekor) Tariklah barisan tak hingga berisi peringkat seragam yang \omterm{def:b2:proba:independence}{saling
bebas}, dalam pengertian kombinatorik berikut: yakni untuk tiap $n$,
urutan relatif $n$ tarikan pertamanya seragam di antara $n!$
kemungkinannya, dan $R_n = {}$``tarikan ke-$n$-nya sebuah rekor (yakni lebih besar
daripada semua sebelumnya)''. Dengan menerima bahwa \omterm{def:b2:proba:space}{kejadian} $R_n$ \omterm{def:b2:proba:independence}{saling
bebas} dengan $\P(R_n) = 1/n$ (buktikanlah sekurangnya kesamaan terakhir
ini lewat kesimetriannya), tunjukkanlah dengan memakai Borel--Cantelli bahwa tak berhingga
banyak rekor terjadi secara hampir pasti, tetapi bahwa rekor pada waktu
yang \emph{berurutan} $n, n+1$ terjadi tak berhingga sering dengan
peluang --- hitunglah $\sum_n \P(R_n \cap R_{n+1})$ lalu
simpulkanlah apa yang diberikan Borel--Cantelli~1.
\end{exercise}

\begin{exercise}[$\star\star\star$]\label{exo:b2:proba:8}
(Bercita rasa Kochen--Stone, versi yang lebih mudah) Misalkan $(A_n)$ \omterm{def:b2:proba:space}{kejadian}
yang \omterm{def:b2:proba:independence}{saling bebas} dengan $\P(A_n) = \frac{1}{n+1}$. Tunjukkan bahwa
$\P(\limsup A_n) = 1$, walaupun $\P(A_n) \to 0$: yakni ``langka
secara individual, pasti secara kolektif''. Sebaliknya, pamerkanlah barisan
\omterm{def:b2:proba:space}{kejadian} (yang bergantungan) dengan $\sum\P(A_n) = \infty$ dan
$\P(\limsup A_n) = 0$, sehingga menunjukkan kebebasannya tak dapat dilepas pada
Borel--Cantelli~2.
\end{exercise}

\begin{exercise}[$\star$]\label{exo:b2:proba:9}
Sekeping koin yang berpeluang gambar $p \in \intoo01$ dilempar sampai
gambar pertamanya. Hitunglah peluang bahwa ini terjadi pada
pangkat yang ganjil, lalu nilailah ia bagi koin yang setimbang.
\end{exercise}

\begin{exercise}[$\star\star$]\label{exo:b2:proba:10}
Misalkan $(A_n)_{n\geq1}$ \omterm{def:b2:proba:independence}{kejadian saling bebas} dengan $\P(A_n) =
p_n < 1$. Tunjukkan bahwa
\[
\P\Bigl(\bigcap_{n\geq1}A_n^c\Bigr)
= \prod_{n\geq1}(1 - p_n)
:= \lim_{N\to\infty}\prod_{n=1}^N(1 - p_n),
\]
dan bahwa limitnya $> 0$ bila dan hanya bila $\sum p_n <
\infty$. Lalu serasikanlah dengan Borel--Cantelli: bahwa ketika $\sum p_n =
\infty$, bukan hanya sebuah $A_n$ terjadi hampir pasti ---
melainkan tak berhingga banyak yang terjadi.
\end{exercise}

\begin{exercise}[$\star\star$]\label{exo:b2:proba:11}
(Kotak korek api Banach) Seorang perokok menyimpan satu kotak berisi $n$ batang korek di
tiap sakunya lalu merogoh saku yang dipilih acak seragam setiap
kalinya. Ketika ia pertama kali menemukan sebuah kotaknya kosong, berapa peluang
bahwa kotak lainnya berisi tepat $k$ batang korek? Tunjukkan bahwa
jawabannya $\binom{2n-k}{n}2^{-(2n-k)}$ lalu periksalah bahwa peluang
itu berjumlah $1$ untuk $n = 1$.
\end{exercise}

\begin{exercise}[$\star\star\star$]\label{exo:b2:proba:12}
(Keaditifan-$\sigma$ itu aksioma yang sejati)
(a) Tunjukkan bahwa tak ada \omterm{def:b2:proba:space}{ukuran peluang} pada $(\N,
\mathcal P(\N))$ yang memberi semua singletonnya bobot yang sama.
(b) Untuk $A \subseteq \N^*$, misalkan $d(A) =
\lim_n\frac{\abs{A\cap\intint1n}}{n}$ ketika limitnya ada
(yakni \emph{kerapatan alaminya}). Tunjukkan bahwa $d$ bersifat aditif berhingga
pada pasangan yang ketiga kerapatannya ada, memberi setiap
singleton kerapatan $0$ dan $\N^*$ kerapatan $1$ --- lalu simpulkan
bahwa $d$ tak bersifat aditif-$\sigma$.
(c) Pamerkanlah sebuah himpunan yang tak berkerapatan. \emph{(Selang-selingkan blok
$\intint{2^{2k}}{2^{2k+1}-1}$ masuk dan keluar.)}
\end{exercise}

\section{Soal: jalan acak sederhana pada $\Z$ bersifat rekuren}

\begin{omfigure}
\begin{tikzpicture}[scale=0.42, y=0.9cm]
  \draw[->] (0,0) -- (25,0) node[below] {$n$};
  \draw[->] (0,-2.6) -- (0,2.9) node[left] {$S_n$};
  \draw[omDef, thick] plot coordinates {(0,0) (1,1) (2,0)
    (3,-1) (4,0) (5,1) (6,2) (7,1) (8,2) (9,1) (10,0) (11,-1)
    (12,-2) (13,-1) (14,0) (15,-1) (16,0) (17,1) (18,2) (19,1)
    (20,0) (21,1) (22,0) (23,1) (24,2)};
  \foreach \t in {2,4,10,14,16,20,22}
    \fill[omThm] (\t,0) circle (5pt);
  \fill (0,0) circle (5pt);
\end{tikzpicture}

{\small Dua puluh empat langkah sebuah \omterm{pb:b2:proba:1}{jalan acak sederhana}; adapun titik
merahnya menandai kepulangan ke titik asalnya. Soal ini menunjukkan bahwa,
dengan peluang $1$, titik itu tak pernah berhenti muncul --- namun
waktu tunggu di antaranya bererata divergen.}
\end{omfigure}

\begin{problem}[{Soal akhir pekan --- teorema kerekurenan Pólya
pada $\Z$, dengan masalah pemungutan suara dan cita rasa arcsinus
di sepanjang jalan}]\label{pb:b2:proba:1}
Lemparkan koin setimbang selamanya; misalkan $X_i = \pm1$ langkah ke-$i$-nya
dan $S_n = X_1 + \dots + X_n$ \emph{jalan acak
sederhana}\index{jalan acak} pada $\Z$, dengan $S_0 = 0$. Seperti pada
\cref{ex:b2:proba:runs}, semua \omterm{def:b2:proba:space}{kejadian} di bawah ditentukan oleh
berhingga banyak lemparan atau berupa kombinasi \omterm{def:b2:structures:countable}{terbilang} \omterm{def:b2:proba:space}{kejadian}
semacam itu, dan kebebasan \omterm{def:b2:proba:space}{kejadian} yang bergantung pada blok lemparan
yang saling lepas merupakan bagian modelnya. Kita menulis $u_n =
\P(S_{2n} = 0)$ dan $N_n(k)$ bagi banyaknya lintasan-$\pm1$
berpanjang $n$ dari $0$ ke $k$.

\medskip
\noindent\textbf{Bagian I --- Mencacah lintasan.}
\begin{enumerate}
  \item Tunjukkan bahwa $N_n(k) = \binom{n}{(n+k)/2}$ ketika $n + k$
        genap dan $\abs k \leq n$, dan $0$ selain itu; lalu turunkan
        $\P(S_n = k) = N_n(k)\,2^{-n}$. Mengapa setiap
        lintasan individual berpanjang $n$ sama mungkinnya?
  \item Tunjukkan $S_{2n+1} \neq 0$ dan $u_n =
        \binom{2n}{n}4^{-n}$, lalu hitunglah $u_1, u_2, u_3$.
  \item Buktikan $u_n = \frac{2n-1}{2n}\,u_{n-1}$; lalu turunkan bahwa
        $(u_n)$ turun ke $0$, dan dari
        \cref{ex:b2:comparison:centralbinomial} bahwa
        \[
        u_n \sim \frac{1}{\sqrt{\pi n}},
        \qquad\text{jadi}\qquad
        \sum_n u_n = \infty .
        \]
  \item (Asas pemantulan) Untuk $k \geq 1$, tunjukkan bahwa
        lintasan berpanjang $n$ dari $1$ ke $k$ yang menyentuh $0$
        berkorespondensi bijektif dengan lintasan dari $-1$ ke $k$;
        lalu turunkan bahwa banyaknya lintasan dari $0$ ke $k$ yang
        tinggal $> 0$ sesudah waktu $0$ adalah $N_{n-1}(k-1) -
        N_{n-1}(k+1)$.
  \item (Teorema pemungutan suara) Turunkan bahwa
        \[
        \P\bigl(S_1 > 0, \dots, S_{n-1} > 0 \bigm| S_n =
        k\bigr) = \frac kn \qquad (k \geq 1) :
        \]
        jadi pada sebuah penghitungan tempat pemenangnya unggul $k$ dari $n$
        surat suara, peluang bahwa pemenangnya unggul sepanjang
        penghitungannya adalah $k/n$. Periksalah dengan tangan untuk $n = 3$ dan $k =
        1$.
\end{enumerate}

\medskip
\noindent\textbf{Bagian II --- Kepulangan ke titik asal.}
\begin{enumerate}[resume]
  \item Buktikan kesamaan kuncinya
        \[
        \P(S_1 \neq 0,\ S_2 \neq 0,\ \dots,\ S_{2n} \neq 0) =
        u_n
        \]
        \emph{(syaratkan pada langkah pertamanya, jumlahkan cacah pada
        pertanyaan 4 atas titik ujungnya, lalu teleskopkan; kemudian rampungkan
        dengan $2\binom{2n-1}{n} = \binom{2n}{n}$)}.
  \item Turunkan dari kekontinuan monoton
        (\cref{thm:b2:proba:continuity}) bahwa jalannya
        pulang ke $0$ sekurangnya sekali dengan peluang $1$,
        dan bahwa $f_n := \P(\text{kepulangan pertama pada waktu }2n)$
        memenuhi
        \[
        f_n = u_{n-1} - u_n = \frac{u_n}{2n-1},
        \qquad \sum_{n\geq1}f_n = 1 .
        \]
  \item Tunjukkan bahwa $\sum_n 2n\,f_n = \infty$: jadi kepulangannya
        pasti, tetapi deret yang akan menghitung waktu tunggu
        reratanya divergen (dalam kosakata
        \cref{ch:b2:randomvar}, waktu kepulangannya
        bernilai harapan tak hingga).
  \item Buktikan bahwa untuk setiap $k \geq 1$ berlaku $\P(\text{sekurangnya }
        k\text{ kepulangan ke }0) = 1$ \emph{(uraikanlah atas
        waktu $k$ kepulangan pertamanya: sebab blok lemparan yang
        bersesuaian saling lepas, jadi peluangnya
        mengalikan lalu berjumlah $(\sum_nf_n)^k$)}; lalu simpulkan dengan
        kekontinuan monoton:
        \[
        \P(S_n = 0 \text{ untuk tak berhingga banyak } n) = 1 :
        \]
        jadi \omterm{pb:b2:proba:1}{jalan acak sederhana} pada $\Z$ bersifat \emph{rekuren}.
  \item Tunjukkan bahwa jalannya mengunjungi setiap tapak $k \in \Z$
        secara hampir pasti, sehingga (berkat kerekurenannya, yang dimulai ulang di
        kunjungan pertamanya) tak berhingga sering. \emph{(Sebab tanda
        ekskursi berurutan dari $0$ merupakan koin setimbang yang \omterm{def:b2:proba:independence}{saling
        bebas}; dan sebuah ekskursi positif mengunjungi $1$.)}
\end{enumerate}

\medskip
\noindent\textbf{Bagian III --- Borel--Cantelli dan jalan yang
berat sebelah.}
\begin{enumerate}[resume]
  \item \omterm{def:b2:proba:space}{Kejadian} $A_n = \{S_{2n} = 0\}$ memenuhi $\sum\P(A_n)
        = \infty$; jelaskanlah mengapa Borel--Cantelli~2 \emph{tak}
        berlaku padanya, dan apa yang akan diberikan Borel--Cantelli~1
        seandainya deretnya konvergen. (Inilah strategi seluruh Bagiannya.)
  \item Kini biarkanlah koinnya berat sebelah dengan $p \neq \frac12$ dan $q = 1 -
        p$. Tunjukkan $\P(S_{2n} = 0) = \binom{2n}n(pq)^n =
        u_n\,(4pq)^n$ dengan $4pq < 1$, lalu turunkan
        $\sum_n\P(S_{2n} = 0) < \infty$, kemudian simpulkan lewat
        Borel--Cantelli~1 bahwa jalan berat sebelahnya pulang ke $0$
        hanya berhingga kali, secara hampir pasti.
  \item Masih untuk $p \neq \frac12$: tunjukkan $\P(S_n = k) \leq
        \binom{n}{\floor{n/2}}\,(pq)^{n/2}\,(p/q)^{k/2}$ untuk
        tiap $k$ yang tetap, lalu turunkan bahwa setiap tapaknya dikunjungi
        berhingga kali secara hampir pasti, kemudian simpulkan
        $\abs{S_n} \to \infty$ secara hampir pasti: jadi jalan berat sebelahnya
        bersifat \emph{transien}.
  \item Kembali ke koin setimbangnya: dengan memakai pertanyaan 6, hitunglah
        peluang bahwa $200$ lemparan menghasilkan \emph{tak ada} seri
        (yakni $S_n \neq 0$ untuk $1 \leq n \leq 200$), yang secara numerik
        $u_{100} \approx 0.056$. Berilah komentar tentang peluruhan
        $1/\sqrt{\pi n}$ yang lambat: bahwa seri itu pasti dalam jangka
        \omterm{def:b2:curves:length}{panjang} tetapi lebih langka daripada dugaan gerak hati.
  \item (Singgahan pertama) Misalkan $T_1$ waktu pertama jalannya
        mengenai $1$. Dengan memakai asas pemantulan bagi maksimumnya
        $M_n = \max_{i\leq n}S_i$ (yang dibuktikan pada pertanyaan
        16, yang tak bergantung pada pertanyaan ini), atau langsung
        dari pertanyaan 7 dengan menyaratkan pada langkah pertamanya,
        tunjukkanlah $\P(T_1 = 2n - 1) = f_n$; lalu turunkan $\P(T_1 <
        \infty) = 1$ sedangkan deret waktu reratanya
        $\sum(2n-1)f_n$ divergen.
\end{enumerate}

\medskip
\noindent\textbf{Bagian IV --- Maksimum, nol terakhir, unggul \omterm{def:b2:curves:length}{panjang}.}
\begin{enumerate}[resume]
  \item (Pemantulan bagi maksimumnya) Untuk $k \geq 1$, buktikanlah
        \[
        \P(M_n \geq k) = 2\,\P(S_n > k) + \P(S_n = k)
        \]
        dengan memantulkan lintasannya sesudah kunjungan pertamanya ke
        aras $k$.
  \item Turunkan $\P(M_{2n} \geq 1) = 1 - u_n$, yakni
        $\P(S_i \leq 0 \text{ untuk setiap } i \leq 2n) = u_n$:
        jadi peluang tak pernah unggul sama dengan
        peluang tak pernah berada di nol (pertanyaan 6) ---
        yakni dua \omterm{def:b2:proba:space}{kejadian} yang berbeda, satu peluang.
  \item (Nol terakhir) Misalkan $L_{2n} = \max\{k \leq 2n : S_k =
        0\}$ (yang genap). Dengan menggabungkan pertanyaan 6 dan kebebasan
        blok lemparan yang saling lepas, tunjukkanlah
        \[
        \P(L_{2n} = 2k) = u_k\,u_{n-k}
        \qquad (0 \leq k \leq n),
        \]
        lalu turunkan, tanpa perhitungan lebih lanjut, kesamaan
        binomial $\sum_{k=0}^n u_ku_{n-k} = 1$.
  \item Tunjukkan bahwa hukum $L_{2n}$ bersifat setangkup (yakni $\P(L =
        2k) = \P(L = 2n - 2k)$) dan, dengan memakai $u_j \sim
        1/\sqrt{\pi j}$, bahwa ekstremnya merupakan nilainya yang paling
        mungkin. Tabelkanlah untuk $n = 5$: $\P(L_{10} = 0)
        = u_5 \approx 0.246$ lawan $\P(L_{10} = 4) = u_2u_3
        \approx 0.117$. Tafsirkanlah: bahwa pada permainan setimbang yang \omterm{def:b2:curves:length}{panjang},
        seri terakhirnya cenderung sangat awal atau sangat lambat ---
        jadi keunggulan yang \omterm{def:b2:curves:length}{panjang} adalah aturannya, bukan kekecualiannya.
  \item Rangkailah pertanyaan 16--19 menjadi sebuah paragraf tentang
        gambaran fluktuasi jalan setimbangnya: yakni skala difusif
        yang disarankan pertanyaan 3, kepastian
        kepulangannya lawan waktu tunggu reratanya yang divergen, dan
        keawetan keunggulannya yang bercita rasa arcsinus.
\end{enumerate}

\medskip
\noindent\textbf{Bagian V --- Kesamaan pembaruan dan
teorema Pólya.}
\begin{enumerate}[resume]
  \item Buktikan, dengan memartisi $\{S_{2n} = 0\}$ atas waktu
        kepulangan pertamanya, \emph{kesamaan pembaruannya}
        \[
        u_n = \sum_{k=1}^n f_k\,u_{n-k} \quad (n \geq 1),
        \qquad\text{sehingga}\qquad
        U(x)\bigl(1 - F(x)\bigr) = 1 \quad (0 \leq x < 1),
        \]
        dengan $U(x) = \sum_{n\geq0}u_nx^n$ dan $F(x) =
        \sum_{n\geq1}f_nx^n$ (lalu benarkanlah jari-jarinya beserta
        hasil kali deretnya lewat \cref{ch:b2:powerseries}).
  \item Turunkan \emph{dikotomi kerekurenannya}: yakni dengan melewatkan $x
        \to 1^-$ (lewat limit monoton deret berkoefisien taknegatif),
        \[
        \sum_n u_n = \infty \iff \sum_n f_n = 1 ,
        \]
        lalu periksalah ia terhadap pertanyaan 3, 7 (jalan setimbangnya) dan
        12 (jalan berat sebelahnya).
  \item (Dimensi $2$) Jalan sederhana pada $\Z^2$ melangkah
        $(\pm1, 0)$ dan $(0, \pm1)$ secara seragam. Tunjukkan bahwa
        koordinat terputarnya $U_n = X_n + Y_n$ dan $V_n = X_n
        - Y_n$ menjalankan jalan setimbang yang \emph{\omterm{def:b2:proba:independence}{saling bebas}} pada
        $\Z$, lalu turunkan
        \[
        \P\bigl(S^{(2)}_{2n} = (0,0)\bigr) = u_n^2 \sim
        \frac1{\pi n},
        \qquad \sum_n u_n^2 = \infty ,
        \]
        kemudian simpulkan dengan pertanyaan 21--22 (yang buktinya
        berpindah kata demi kata) bahwa jalannya pada $\Z^2$ bersifat
        rekuren.
  \item (Dimensi $3$) Untuk jalan sederhana pada $\Z^3$, terimalah
        taksiran lokalnya $\P(S^{(3)}_{2n} = 0) \leq
        C\,n^{-3/2}$ (yang dibuktikan lewat teorema limit lokal pada
        jilid Tahun ke-3). Turunkan dari Borel--Cantelli~1
        bahwa jalannya pada $\Z^3$ bersifat transien, lalu nyatakanlah
        hasil penuhnya: yakni \emph{teorema Pólya} --- bahwa \omterm{pb:b2:proba:1}{jalan
        acak sederhana} bersifat rekuren pada dimensi $1$ dan $2$, dan
        transien pada dimensi $3$ ke atas.
  \item Rangkuman. Daftarkanlah peran persis yang dimainkan: pencacahan
        lintasan dan pemantulan; kekontinuan monoton;
        kebebasan blok lemparan yang saling lepas;
        Borel--Cantelli~1; dan kesamaan pembaruannya. Lalu fakta
        analitis tunggal mana ($u_n \sim 1/\sqrt{\pi n}$, sehingga
        $\sum u_n = \infty$ tetapi $\sum u_n^2 = \infty$ dan
        $\sum n^{-3/2} < \infty$) yang memutuskan antara kerekurenan
        dan ketransienan pada tiap dimensinya?

\end{enumerate}
\end{problem}
